% Source: https://artofproblemsolving.com/community/c3h1368794_mathcounts_graph_geometry
% Stable identity: aops-topic-1368794-q01
The slope is
\[m=\dfrac{-13-8}{5-(-2)}=\dfrac{-21}{7}=-3 .\]
Substituting $(-2,8)$ into $y=-3x+b$ gives $8=6+b$, so $b=2$. (Check with the other point: $-3\cdot5+2=-13$.)

\begin{asy}
usepackage("amsmath");
size(4.5cm,6.5cm);
defaultpen(fontsize(10pt));
import graph;
real f(real x){return -3x+2;}
draw((-3.4,0)--(6,0),gray,Arrow(4)); draw((0,-15)--(0,10),gray,Arrow(4));
draw(graph(f,-2.6,5.3),blue+linewidth(0.9));
dot((-2,8)); dot((5,-13)); dot((0,2),red);
label("$(-2,8)$",(-2,8),W); label("$(5,-13)$",(5,-13),E); label("$(0,2)$",(0,2),E,red);
label("$x$",(6,0),SE); label("$y$",(0,10),NE);
label("\mbox{}",(min(currentpicture,user=true)+max(currentpicture,user=true))/2,black); // keeps the text after the figure black when LaTeX typesets the labels inline
\end{asy}

So the line is $y=-3x+2$ and $mb=(-3)(2)=\boxed{-6}$.
